% Рисунок к примеру: функция Хевисайда и f(x) = x^2 sign x.
\begin{figure}[htbp]
\centering
\begin{tikzpicture}
\begin{groupplot}[
  group style={group size=2 by 1, horizontal sep=1.8cm},
  width=7.2cm, height=4.2cm,
  axis lines=middle,
  samples=120, clip=false,
  xtick=\empty,
]
% --- Хевисайд
\nextgroupplot[xmin=-1.3, xmax=1.3, ymin=-0.4, ymax=1.5,
               ytick={1}, yticklabels={$1$}]
  \addplot[very thick, blue!60!black] coordinates {(-1.2,0) (0,0)};
  \addplot[very thick, blue!60!black] coordinates {(0,1) (1.2,1)};
  \addplot[only marks, mark=*, blue!60!black] coordinates {(0,1)};
  \addplot[only marks, mark=o, blue!60!black] coordinates {(0,0)};
  \node[anchor=south east, font=\footnotesize] at (axis cs:1.2,1.02) {$H$};
  \node[anchor=north, font=\footnotesize] at (axis cs:0,-0.12) {$H' = \delta$};

% --- f(x) = x^2 sign x
\nextgroupplot[xmin=-1.3, xmax=1.3, ymin=-1.7, ymax=1.7, ytick=\empty]
  \addplot[very thick, blue!60!black, domain=0:1.2] {x^2};
  \addplot[very thick, blue!60!black, domain=-1.2:0] {-x^2};
  \node[blue!60!black, anchor=south east, font=\footnotesize] at (axis cs:1.25,1.3)
       {$f = x^2 \sgn x$};
  \node[anchor=north, font=\footnotesize] at (axis cs:0,-1.05) {$f''' = 4\delta$};
\end{groupplot}
\end{tikzpicture}
\caption{Слева: у функции Хевисайда производная в смысле распределений равна $\delta$,
а $\delta$ сингулярна — значит, соболевской производной у $H$ нет. Справа:
$f = x^2\sgn x$ дифференцируема по Соболеву дважды ($f' = 2|x|$, $f'' = 2\sgn x$),
но третья производная уже сингулярна.}
\label{fig:heaviside}
\end{figure}
